\documentclass[10pt]{amsart}

\usepackage[margin=1in]{geometry}
\usepackage[T1]{fontenc}
\usepackage{lmodern}
\usepackage{amsmath,amssymb,amsthm,mathtools}
\usepackage{microtype}
\usepackage{url}
\usepackage{enumitem}
\usepackage{alltt}
\usepackage{booktabs}
\usepackage[dvipsnames]{xcolor}
\usepackage[colorlinks=true,linkcolor=MidnightBlue,citecolor=OliveGreen,urlcolor=MidnightBlue]{hyperref}
\usepackage[capitalize,nameinlink]{cleveref}

\setlist{leftmargin=*,itemsep=2pt,topsep=4pt}
\emergencystretch=2.5em

\theoremstyle{plain}
\newtheorem{theorem}{Theorem}[section]
\newtheorem{lemma}[theorem]{Lemma}
\newtheorem{proposition}[theorem]{Proposition}
\newtheorem{corollary}[theorem]{Corollary}
\theoremstyle{definition}
\newtheorem{definition}[theorem]{Definition}
\newtheorem{example}[theorem]{Example}
\theoremstyle{remark}
\newtheorem{remark}[theorem]{Remark}
\numberwithin{equation}{section}

\newcommand{\A}{\textup{[A]}}
\newcommand{\N}{\mathbb{N}}
\newcommand{\EC}{\operatorname{ec}}
\newcommand{\HC}{\operatorname{hc}}
% Lean names: typewriter, with line breaks allowed after '.' and '_'.
\DeclareUrlCommand\lean{\urlstyle{tt}\def\UrlBreaks{\do\.\do\_}\def\UrlBigBreaks{}}
% Repository names: typewriter, with line breaks allowed after '/' and '-'.
\DeclareUrlCommand\repo{\urlstyle{tt}\def\UrlBreaks{\do\/\do\-}\def\UrlBigBreaks{}}

\begin{document}

\title[Counting Eulerian circuits in subexponential time]{Counting Eulerian circuits in subexponential time,\\ assuming the new $\#$SAT algorithms}

\author{Vaibhav Tulsyan}

\date{November 2026}

\subjclass[2020]{05C45, 05C30, 68Q25, 68W40, 03B35}
\keywords{Eulerian circuits, exact counting, transition systems, line graphs, Hamiltonian cycles, Exponential Time Hypothesis, Lean~4, formal verification}

\begin{abstract}
Counting the Eulerian circuits of an undirected graph is $\#$P-complete, even for planar $4$-regular
graphs, while for directed graphs the BEST theorem counts them in polynomial time. However, we are
unaware of any exact algorithm for undirected graphs of bounded degree that runs in time $2^{o(m)}$,
where $m$ is the number of edges; the simple methods enumerate transition systems or Eulerian
orientations. We show that such an algorithm follows from the exact counting of Hamiltonian cycles
that we proved earlier on top of Anthropic's claimed $\#$SAT algorithm. The reduction is classical:
by Kotzig's description, the Eulerian circuits of $G$ are in bijection with the Hamiltonian cycles of
the line graph of the double subdivision of $G$, which has $3m$ vertices and at most $(\Delta+1)m$
edges. We prove in Lean~4 (\texttt{theoremEuler}, \texttt{theoremEuler\_noIso},
\texttt{theoremEuler\_rate}; axioms \texttt{propext}, \texttt{Classical.choice}, \texttt{Quot.sound})
that for every $\delta>0$ and $\Delta$ a finite-alphabet Mathlib Turing machine outputs the number of
Eulerian circuits of every graph of maximum degree at most $\Delta$ within
$\lfloor C(n+2)^C2^{\delta m}\rfloor$ steps, and within $\lfloor C2^{\delta m}\rfloor$ steps when there is
no isolated vertex; and that one machine does so for all graphs within
$\lfloor C(n+2)^C2^{Cm\log_2(\Delta+2)/\sqrt{\log_2(m+2)}}\rfloor$ steps. The factor $(n+2)^C$ pays for
reading the $n^2$-bit adjacency matrix and cannot be dropped. The bijection is checked by brute force
on twelve small graphs, for example $132$ circuits for $K_5$. The formal object is Kotzig's
transition-system definition; its agreement with Mathlib's Eulerian walks is not formalised. The
constants are astronomical, and the proof rests on Anthropic's Lean development of its claimed
refutation of the Exponential Time Hypothesis, which is under review.
\end{abstract}

\maketitle

% =====================================================================
\section{Introduction}\label{sec:intro}
% =====================================================================

An \emph{Eulerian circuit} of a graph is a closed trail that uses every edge exactly once. Deciding
whether one exists is easy (Euler: the graph is connected apart from isolated vertices and every
degree is even), but counting them is not. For directed graphs the BEST theorem of de Bruijn and van
Aardenne-Ehrenfest~\cite{BE51}, building on Smith and Tutte~\cite{ST41}, expresses the number as a
spanning-tree count times $\prod_v(d^+_v-1)!$, so it is computable in polynomial time. For undirected
graphs Brightwell and Winkler~\cite{BW05} proved the problem $\#$P-complete, and Ge and
\v{S}tefankovi\v{c}~\cite{GS10} proved it $\#$P-complete even for planar $4$-regular graphs.

However, the exact algorithms we know for general undirected graphs take time exponential in the
number $m$ of edges, even when the degree is bounded. The simplest one lists the
$\prod_v(d_v-1)!!$ \emph{transition systems} (a pairing of the edge ends at every vertex) and keeps
those that form one closed trail. Another sums the BEST formula over the Eulerian orientations,
of which a $2k$-regular graph has at least $(\binom{2k}{k}/2^k)^n$ by Schrijver~\cite{Sch83}, so also
exponentially many in $m$. Polynomial algorithms exist for restricted classes: generalised
series-parallel graphs~\cite{CCM12} and graphs of bounded treewidth~\cite{BDG15}. We are unaware of
an exact algorithm that runs in time $2^{o(m)}$ on all graphs of bounded degree, even conditionally.

\subsection{The problem and the running example}
We consider the problem $\#\mathrm{EC}$: given a finite simple graph $G$ with vertex set
$\{0,\dots,n-1\}$ as its adjacency matrix, output the number $\EC(G)$ of its Eulerian circuits,
each counted once up to the choice of starting point and direction. Following Kotzig~\cite{Kot68}, we
\emph{define} this number through transitions (\cref{def:kotzig}): $\EC(G)$ is the number of
transition systems whose trails form a single closed trail. With this definition $\EC(G)=0$ if some
degree is odd, if the edges do not lie in one component, or if $G$ has no edges.

\begin{example}[Running example]\label{ex:k5}
Take $G=K_5$, so $n=5$, $m=10$ and every degree is $4$. At each vertex the four edge ends can be
paired in $3$ ways, so there are $3^5=243$ transition systems. Splitting them by the number of closed
trails they form gives $132$ with one trail, $96$ with two and $15$ with three. So $\EC(K_5)=132$.
(OEIS A007082~\cite{OEIS} lists $264$ for $K_5$, counting the two directions separately, and records
that $132$ tours start with a fixed directed edge.)
\end{example}

\subsection{The naive route and the principle}
At first one might encode $\#\mathrm{EC}$ directly as a counting problem for a $\#$SAT engine:
a variable block per vertex for its transition, and a constraint that the trails form one circuit.
However, ``one circuit'' is a global connectivity constraint, and the engines of
OpenAI~\cite{OpenAI26} and Anthropic~\cite{Anthropic26} (below \A) count only local constraints;
encoding connectivity was the whole difficulty of our earlier work on Hamiltonian
cycles~\cite{Tulsyan26}. Instead, we rely on a different principle: \emph{connectivity has already
been paid for once, in the Hamiltonian-cycle counter, so a problem whose global constraint is
Hamiltonicity of a sparse graph needs only a polynomial-time reduction.} Kotzig's description provides
exactly that reduction. Subdivide every edge of $G$ twice and take the line graph. The result $G'$
has $3m$ vertices, and its Hamiltonian cycles are the transition systems of $G$ with one trail
(\cref{sec:bijection}). For $K_5$, $G'$ has $30$ vertices and $50$ edges, and $132$ Hamiltonian
cycles.

\subsection{What we prove}
In this paper we compose this reduction with our machine-checked Hamiltonian-cycle
counters~\cite{Tulsyan26}, and prove the result in Lean. Write $\Delta$ for the maximum degree of
$G$. Informally, the three theorems of \cref{sec:theorems} say the following.

\begin{theorem}[informal; machine-checked]\label{thm:informal}
\leavevmode
\begin{enumerate}[label=\textup{(\alph*)}]
\item For every $\delta>0$ and $\Delta$ there are a Turing machine and a constant $C$ that output
$\EC(G)$ for every $G$ of maximum degree at most $\Delta$ within $C(n+2)^C2^{\delta m}$ steps.
\item The same holds within $C2^{\delta m}$ steps for graphs without isolated vertices.
\item One machine and one $C$ output $\EC(G)$ for every $G$ within
$C(n+2)^C\,2^{Cm\log_2(\Delta+2)/\sqrt{\log_2(m+2)}}$ steps.
\end{enumerate}
\end{theorem}

\noindent Part~(c) is subexponential in $m$ whenever $\log\Delta=o(\sqrt{\log m})$; for bounded
degree it is $2^{O(m/\sqrt{\log m})}$.

\subsection{Contributions}
Our contributions are:
\begin{itemize}
\item We state the problem $\#\mathrm{EC}$ in Kotzig's form as Lean definitions short enough to check
by eye, and list them (\cref{sec:defs}).
\item We prove in Lean that the Hamiltonian cycles of the line graph of the double subdivision are in
bijection with the Eulerian circuits, and bound its size by $3m$ vertices and $2(\Delta+1)m$ edges
(\cref{sec:bijection}).
\item We prove three Turing-machine theorems for $\#\mathrm{EC}$, with bounds $2^{\delta m}$ (bounded
degree) and $2^{O(m\log\Delta/\sqrt{\log m})}$ (one machine), and explain why the factor $(n+2)^C$ is
needed (\cref{sec:theorems}).
\item We check the bijection numerically on twelve graphs in three independent ways
(\cref{sec:numerics}), and describe the Lean development and its trusted base (\cref{sec:formal}).
\end{itemize}
The bijection is classical; Balaji, Datta and Ganesan~\cite{BDG15} count Euler tours of a graph as the
Hamiltonian cycles of the line graph of a modified graph. What we add is the
subexponential count, conditional on \A, and its machine-checked proof. To our knowledge no
$2^{o(m)}$ exact algorithm for bounded-degree graphs was stated before, even conditionally; we make
no claim of practicality.

\subsection{Status of each claim}
We separate four kinds of claim.
\emph{Machine-checked in Lean}: \cref{thm:euler,thm:noiso,thm:rate}, as
\lean{EthFollowup.ExactHC.theoremEuler}, \lean{EthFollowup.ExactHC.theoremEuler_noIso} and
\lean{EthFollowup.ExactHC.theoremEuler_rate}, and \cref{lem:bij,lem:size}, as
\lean{ESplit.card_hamCycles_eulerSplit} and \lean{ESplit.card_edgeFinset_eulerSplit_le}; all with the
axioms \lean{propext}, \lean{Classical.choice} and \lean{Quot.sound} only.
\emph{Proved on paper only}: the agreement of Kotzig's definition with closed trails up to rotation
and reversal (\cref{rem:kotzig}), the lower bound behind the factor $(n+2)^C$ (\cref{prop:lower}), and
the reduction for multigraphs (\cref{rem:multi}).
\emph{Assumed}: nothing beyond the three axioms. \A's results and our Hamiltonian-cycle theorems
enter as Lean theorems, which the kernel checks with ours; but the result stands or falls with \A's
development (\cref{sec:limitations}).
\emph{Checked numerically}: the table of \cref{sec:numerics} and the counts in \cref{ex:k5}.

\subsection*{Outline}
\Cref{sec:defs} gives the definitions, \cref{sec:bijection} the bijection, \cref{sec:theorems} the
theorems, \cref{sec:numerics} the numerical checks, \cref{sec:formal} the Lean development, and
\cref{sec:related,sec:limitations} related work and limitations.

% =====================================================================
\section{Kotzig's definition, as a reader must check it}\label{sec:defs}
% =====================================================================

In this section we give the definitions on paper and then in Lean. Everything a reader must check to
trust the theorems of \cref{sec:theorems} is here, together with Mathlib and \A's input format
(\cref{sec:theorems}).

Let $G$ be a simple graph on $[n]=\{0,\dots,n-1\}$ and list its edges $e_i=\{u_i,v_i\}$, $u_i<v_i$,
in lexicographic order, $i<m$. Each edge has two \emph{ends}, its end at $u_i$ and its end at $v_i$.

\begin{definition}[Kotzig]\label{def:kotzig}
A \emph{transition system} of $G$ is a choice, at every vertex $x$, of a perfect matching of the ends
at $x$; a matched pair is a \emph{transition}. Following the transitions, the edges fall into closed
trails. An \emph{Eulerian circuit} is a transition system with exactly one closed trail, and $\EC(G)$
is their number.
\end{definition}

\begin{remark}[on paper]\label{rem:kotzig}
For $m\ge1$, a closed trail through every edge, taken up to rotation and reversal, determines its
transition system (consecutive edges at each passage through a vertex), and a transition system with
one trail determines the circuit up to rotation and reversal. This is the standard
correspondence~\cite{Kot68}; we use it only to explain the name, and we do not formalise it.
\end{remark}

\paragraph{The split graph.} The graph $G'$ has vertex set $[3m]$. Edge $e_i$ owns three
\emph{copies}: $\mathrm{lo}_i=3i$ (its end at $u_i$), $\mathrm{mid}_i=3i+1$, and $\mathrm{hi}_i=3i+2$
(its end at $v_i$). Its edges are the \emph{internal} edges $\{\mathrm{lo}_i,\mathrm{mid}_i\}$,
$\{\mathrm{mid}_i,\mathrm{hi}_i\}$, and, for every vertex $x$ and distinct edges $e_i,e_j$ at $x$, the
\emph{transition edge} between the copies of $e_i$ and $e_j$ at $x$. So the copies at $x$ form a
clique. Subdividing $e_i$ as $u_i\,a_i\,b_i\,v_i$ and naming the three pieces
$\mathrm{lo}_i=u_ia_i$, $\mathrm{mid}_i=a_ib_i$, $\mathrm{hi}_i=b_iv_i$ shows that $G'$ is the line
graph of the double subdivision of $G$. A transition system of $G$ is exactly a set $T$ of transition
edges with one edge at every $\mathrm{lo}$ and every $\mathrm{hi}$ copy, and its \emph{transition
graph} is $T$ together with the internal edges.

\paragraph{In Lean.} The definitions are in \lean{lean/EthFollowup/ExactHC/EulerSplit.lean}, namespace
\lean{EthFollowup.ExactHC.ESplit}, with \texttt{variable \{n : $\mathbb{N}$\} (G : SimpleGraph (Fin
n)) [DecidableRel G.Adj]}. The transition predicate \texttt{A x y z} allows or forbids the
transition at $x$ between the edges $xy$ and $xz$; the theorems use \lean{eulerSplitAll} and
\lean{eulerCircuitsAll}, where every transition is allowed. We copy the definitions, rendering Unicode
symbols as mathematics and omitting docstrings and the proof terms in \lean{cLo}, \lean{cMid},
\lean{cHi}:
\begin{alltt}\small
def edgeList : List (Fin n \(\times\) Fin n) :=
  (List.finRange n).flatMap fun u =>
    ((List.finRange n).filter fun v => decide (u < v \(\wedge\) G.Adj u v)).map fun v => (u, v)
abbrev M : \(\mathbb{N}\) := (edgeList G).length
def eLo (i : Fin (M G)) : Fin n := ((edgeList G).get i).1
def eHi (i : Fin (M G)) : Fin n := ((edgeList G).get i).2
def cLo (i : Fin (M G)) : Fin (3 * M G) := \(\langle\)3 * i.val, _\(\rangle\)      -- cMid: 3i+1, cHi: 3i+2
def copyAt (i : Fin (M G)) (x : Fin n) : Fin (3 * M G) :=
  if eLo G i = x then cLo G i else cHi G i
def other (i : Fin (M G)) (x : Fin n) : Fin n := if eLo G i = x then eHi G i else eLo G i
def intE : Finset (Sym2 (Fin (3 * M G))) :=
  univ.image (fun i => s(cLo G i, cMid G i)) \(\cup\) univ.image (fun i => s(cMid G i, cHi G i))
def cliqueE (A : Fin n \(\to\) Fin n \(\to\) Fin n \(\to\) Bool) : Finset (Sym2 (Fin (3 * M G))) :=
  (univ.filter fun t : Fin (M G) \(\times\) Fin (M G) \(\times\) Fin n => t.1 \(\ne\) t.2.1 \(\wedge\)
      (eLo G t.1 = t.2.2 \(\vee\) eHi G t.1 = t.2.2) \(\wedge\) (eLo G t.2.1 = t.2.2 \(\vee\) eHi G t.2.1 = t.2.2) \(\wedge\)
      A t.2.2 (other G t.1 t.2.2) (other G t.2.1 t.2.2) = true).image
    fun t => s(copyAt G t.1 t.2.2, copyAt G t.2.1 t.2.2)
abbrev eulerSplit (A) : SimpleGraph (Fin (3 * M G)) := graphOf (intE G \(\cup\) cliqueE G A)
abbrev eulerSplitAll : SimpleGraph (Fin (3 * M G)) := eulerSplit G fun _ _ _ => true
def transSys (A) : Finset (Finset (Sym2 (Fin (3 * M G)))) :=
  (cliqueE G A).powerset.filter fun T => \(\forall\) i : Fin (M G),
    (T.filter fun e => cLo G i \(\in\) e).card = 1 \(\wedge\) (T.filter fun e => cHi G i \(\in\) e).card = 1
def phiE (T : Finset (Sym2 (Fin (3 * M G)))) : Finset (Sym2 (Fin (3 * M G))) := T \(\cup\) intE G
def eulerCircuits (A) : Finset (Finset (Sym2 (Fin (3 * M G)))) :=
  (transSys G A).filter fun T => (graphOf (phiE G T)).Connected
abbrev eulerCircuitsAll : Finset (Finset (Sym2 (Fin (3 * M G)))) :=
  eulerCircuits G fun _ _ _ => true
\end{alltt}
Here \texttt{(A)} abbreviates \texttt{(A : Fin n $\to$ Fin n $\to$ Fin n $\to$ Bool)}. The helpers come
from our earlier development~\cite{Tulsyan26} (module \lean{HamiltonianParity}):
\lean{graphOf X} is \lean{SimpleGraph.fromEdgeSet X}; \lean{TwoRegular X} says that every vertex
meets exactly two edges of \texttt{X}; \lean{twoFactors G} is
\lean{G.edgeFinset.powerset.filter TwoRegular}; and \lean{hamCycles G} is the set of those
\texttt{X} in \lean{twoFactors G} with \lean{(graphOf X).Connected}. Connectedness is Mathlib's, so
a graph on zero vertices is not connected, and $\EC(G)=0$ when $m=0$. The trails of $T$ are the
components of the transition graph, so \lean{eulerCircuitsAll G} is \cref{def:kotzig} verbatim.

% =====================================================================
\section{Hamiltonian cycles of the split graph are Eulerian circuits}\label{sec:bijection}
% =====================================================================

In this section we prove the bijection and bound the size of $G'$.

\begin{lemma}[\lean{ESplit.hamCycles_eulerSplit}, \lean{ESplit.card_hamCycles_eulerSplit}; machine-checked]\label{lem:bij}
For every transition predicate $A$,
\[
  \texttt{hamCycles (eulerSplit G A)} = \texttt{(eulerCircuits G A).image (phiE G)},
\]
and \lean{phiE G} is injective on \lean{eulerCircuits G A}. Hence
\[
  \texttt{(hamCycles (eulerSplit G A)).card = (eulerCircuits G A).card},
\]
and in particular $\HC(G')=\EC(G)$.
\end{lemma}

\begin{proof}
A middle copy $\mathrm{mid}_i$ has exactly two neighbours, $\mathrm{lo}_i$ and $\mathrm{hi}_i$. So
every Hamiltonian cycle $H$ of $G'$, indeed every $2$-factor, contains all internal edges, and then
uses exactly one transition edge at every $\mathrm{lo}$ and every $\mathrm{hi}$ copy: the set $T$ of
its transition edges is a transition system, and $H=T\cup\mathrm{intE}$
(\lean{twoRegular_phiE_iff}). Conversely, for every transition system $T$ the set
$T\cup\mathrm{intE}$ is a $2$-factor, and it is a Hamiltonian cycle exactly when it is connected.
Transition edges and internal edges are disjoint, so $T$ is recovered from $H$ as $H\cap\mathrm{cliqueE}$
(\lean{filter_phiE}), which gives injectivity.
\end{proof}

\noindent In the running example a Hamiltonian cycle of $G'$ for $K_5$ walks
$\mathrm{lo}_i\,\mathrm{mid}_i\,\mathrm{hi}_i$ along every edge and jumps between edges through a
transition edge at the shared vertex; the $243-132=111$ transition systems with two or three trails
give the $2$-factors of $G'$ with two or three cycles.

\begin{lemma}[\lean{ESplit.card_edgeFinset_eulerSplit_le}, \lean{Eul.card_edge_eulerSplitAll_le}; machine-checked]\label{lem:size}
$G'$ has $3m$ vertices (\lean{ESplit.length_edgeList}: \lean{M G} $=|E(G)|$) and at most
$2(\Delta+1)m$ edges; if $\Delta\le D$, then $|E(G')|\le(D+1)\cdot3m$.
\end{lemma}

\noindent On paper, $|E(G')|=2m+\sum_x\binom{d_x}{2}\le(\Delta+1)m$, and the maximum degree of $G'$
is $\max(2,\Delta)$ when $m\ge1$: a copy at $x$ has $d_x-1$ transition neighbours and one middle
neighbour. The Lean bound is cruder by a factor~$2$, which is harmless. For $K_5$:
$20+5\binom42=50$.

\paragraph{The preprocessing.} The machine of \cref{sec:theorems} must write the adjacency matrix of
$G'$ from that of $G$. We prove an arithmetic description of the edges: for copies $x\ne y$ of $[3m]$,
$\{x,y\}$ is internal iff $\lfloor x/3\rfloor=\lfloor y/3\rfloor$ and $x\bmod3+y\bmod3\in\{1,3\}$, and is
a transition edge iff $\lfloor x/3\rfloor\ne\lfloor y/3\rfloor$, neither residue is $1$, and the two
copies sit at the same vertex of $G$ (\lean{mem_intE_arith}, \lean{mem_cliqueE_arith},
\lean{adj_eulerSplitAll_iff}). The word map \lean{Eul.eulW} lists the edges from the matrix
(\lean{Eul.edgeListW}), writes $3m$ in unary and the $(3m)^2$ entries, and satisfies
\lean{eulW_code : eulW (code G) = code (eulerSplitAll G)}; it is polynomial-time in \A's sense
(\lean{eulW_polyT}).

% =====================================================================
\section{The theorems}\label{sec:theorems}
% =====================================================================

In this section we state the three theorems exactly as in Lean, explain the factor $(n+2)^C$, and
sketch the proofs. The theorems are in
\lean{lean-bridge-anthropic/bridge/EthFollowup/ExactHC/EulerTheorems.lean}, namespace
\lean{EthFollowup.ExactHC}. They use \A's input format and running-time predicate, from \A's file
\texttt{GraphProblemsStatement.lean}: \lean{GraphProblems.code G} is $n$ in unary ($n$ ones and a
zero) followed by the $n^2$ entries of the adjacency matrix row by row, a word of length $n^2+n+1$;
and \lean{GraphProblems.Outputs M x y T} says that the Mathlib machine \texttt{M}
(\lean{Turing.TM2ComputableAux}) started on \texttt{x} halts within \texttt{T} steps with output
\texttt{y} (Mathlib's \lean{TM2OutputsInTime}). \lean{Nat.bits} is Mathlib's little-endian binary
representation.

\begin{theorem}[\lean{EthFollowup.ExactHC.theoremEuler}; machine-checked]\label{thm:euler}
In Lean:
\begin{alltt}\footnotesize
theorem theoremEuler (\(\delta\) : \(\mathbb{R}\)) (h\(\delta\) : 0 < \(\delta\)) (\(\Delta\) : \(\mathbb{N}\)) :
    \(\exists\) (Mc : Turing.TM2ComputableAux Bool Bool) (C : \(\mathbb{N}\)), (\(\forall\) i, Finite (Mc.tm.\(\Gamma\) i)) \(\wedge\)
      \(\forall\) (n : \(\mathbb{N}\)) (G : SimpleGraph (Fin n)) [DecidableRel G.Adj], G.maxDegree \(\le\) \(\Delta\) \(\to\)
        GraphProblems.Outputs Mc (GraphProblems.code G) (Nat.bits (eulerCircuitsAll G).card)
          \(\lfloor\)(C : \(\mathbb{R}\)) * ((n : \(\mathbb{R}\)) + 2) ^ C *
            (2 : \(\mathbb{R}\)) ^ (\(\delta\) * (G.edgeFinset.card : \(\mathbb{R}\)))\(\rfloor\)\(\sb{+}\)
\end{alltt}
\end{theorem}

\begin{theorem}[\lean{EthFollowup.ExactHC.theoremEuler_noIso}; machine-checked]\label{thm:noiso}
In Lean:
\begin{alltt}\footnotesize
theorem theoremEuler_noIso (\(\delta\) : \(\mathbb{R}\)) (h\(\delta\) : 0 < \(\delta\)) (\(\Delta\) : \(\mathbb{N}\)) :
    \(\exists\) (Mc : Turing.TM2ComputableAux Bool Bool) (C : \(\mathbb{N}\)), (\(\forall\) i, Finite (Mc.tm.\(\Gamma\) i)) \(\wedge\)
      \(\forall\) (n : \(\mathbb{N}\)) (G : SimpleGraph (Fin n)) [DecidableRel G.Adj], G.maxDegree \(\le\) \(\Delta\) \(\to\)
        (\(\forall\) v, 0 < G.degree v) \(\to\)
        GraphProblems.Outputs Mc (GraphProblems.code G) (Nat.bits (eulerCircuitsAll G).card)
          \(\lfloor\)(C : \(\mathbb{R}\)) * (2 : \(\mathbb{R}\)) ^ (\(\delta\) * (G.edgeFinset.card : \(\mathbb{R}\)))\(\rfloor\)\(\sb{+}\)
\end{alltt}
\end{theorem}

\begin{theorem}[\lean{EthFollowup.ExactHC.theoremEuler_rate}; machine-checked]\label{thm:rate}
In Lean:
\begin{alltt}\footnotesize
theorem theoremEuler_rate :
    \(\exists\) (Mc : Turing.TM2ComputableAux Bool Bool) (C : \(\mathbb{N}\)), (\(\forall\) i, Finite (Mc.tm.\(\Gamma\) i)) \(\wedge\)
      \(\forall\) (n : \(\mathbb{N}\)) (G : SimpleGraph (Fin n)) [DecidableRel G.Adj],
        GraphProblems.Outputs Mc (GraphProblems.code G) (Nat.bits (eulerCircuitsAll G).card)
          \(\lfloor\)(C : \(\mathbb{R}\)) * ((n : \(\mathbb{R}\)) + 2) ^ C *
            (2 : \(\mathbb{R}\)) ^ ((C : \(\mathbb{R}\)) * (G.edgeFinset.card : \(\mathbb{R}\)) *
              Real.logb 2 ((G.maxDegree : \(\mathbb{R}\)) + 2) /
                Real.sqrt (Real.logb 2 ((G.edgeFinset.card : \(\mathbb{R}\)) + 2)))\(\rfloor\)\(\sb{+}\)
\end{alltt}
\end{theorem}

\noindent In words: with $m=|E(G)|$, the bounds are $\lfloor C(n+2)^C2^{\delta m}\rfloor$,
$\lfloor C2^{\delta m}\rfloor$ and
$\lfloor C(n+2)^C2^{Cm\log_2(\Delta+2)/\sqrt{\log_2(m+2)}}\rfloor$, where $\lfloor\cdot\rfloor$ is
\lean{Nat.floor}.

\subsection{The factor \texorpdfstring{$(n+2)^C$}{(n+2)\^{}C} is needed}
Our planned statement of \cref{thm:euler} had the bound $C2^{\delta m}$ alone. That statement is
false, for a reason that has nothing to do with Eulerian circuits: the input is the $n^2$-bit
adjacency matrix, and a graph may consist mostly of isolated vertices, so that $n$ grows while $m$ and
the bound stay constant.

\begin{proposition}[on paper]\label{prop:lower}
Let $f\colon\N\to\N$. No machine in Mathlib's \lean{TM2} model outputs \lean{Nat.bits} $\EC(G)$ on
\lean{code G} within $f(|E(G)|)$ steps for all graphs $G$.
\end{proposition}

\begin{proof}
The machine reads its input from the top of a stack, one symbol per step, so within $t$ steps it sees
at most the first $t$ symbols. Let $E_n$ be the edgeless graph on $[n]$ and $E_n^+$ the same graph with
a triangle on $n-3,n-2,n-1$. Their codes agree on the first $n+1+(n-3)n$ symbols, the triangle's
entries lie in the last three rows, and $\EC(E_n)=0\ne1=\EC(E_n^+)$. Both graphs have at most $3$
edges, so a machine within $f(|E|)$ steps halts after at most $\max(f(0),f(3))$ steps on both, sees
the same prefix once $n^2-2n+1>\max(f(0),f(3))$, and gives the same output. That is a contradiction.
\end{proof}

\noindent The same argument shows that $\Omega(n^2)$ steps are needed on $E_n$, so a factor
polynomial in $n$ is the right repair. When $G$ has no isolated vertex, $n\le2m$, so $(n+2)^C\le
(2m+2)^C$ is absorbed by $2^{(\delta/2)m}$ at the price of a larger constant; this is
\cref{thm:noiso}, the planned form. The factor is a matter of the input format, not of the method:
our Hamiltonian-cycle theorems avoid it only because their bounds already grow with $n$.

\subsection{Proofs}
\emph{\Cref{thm:euler}.} We use the step-level form of our fixed-density Hamiltonian-cycle machine,
\lean{hc_stepped} (in \lean{DirectedTheorems}, the stepped machine behind \lean{theorem1}
of~\cite{Tulsyan26}), with parameters $\delta/3$ and density $D=\Delta+1$. Precompose it with
\lean{eulW} (\lean{Stepped.mapIn}; polynomial overhead by \lean{eulW_polyT}). On \lean{code G} the
composed machine runs the Hamiltonian-cycle machine on \lean{code (eulerSplitAll G)}, a graph with
$n'=3m$ vertices and at most $(\Delta+1)n'$ edges (\cref{lem:size}), and so outputs
$\HC(G')=\EC(G)$ (\cref{lem:bij}) within $A\cdot2^{(\delta/6)\cdot3m}$ simulated steps, with a
polynomial overhead in $n'$ and in the input length $n^2+n+1$. \A's host lemma
\lean{CountKSat.exists_tmW} turns the step-level machine into a finite-alphabet Mathlib machine with
overhead $q(T+1)(|x|+S+1)^q$. Collecting terms gives $C(n+2)^C2^{\delta m}$.

\emph{\Cref{thm:noiso}.} Apply \cref{thm:euler} with $\delta/2$, and bound
$(n+2)^C\le(2m+2)^C\le C_p2^{(\delta/4)m}$ (\lean{exists_pow_le}).

\emph{\Cref{thm:rate}.} We use the machine \lean{T3.hcM3} of our maximum-degree theorem
\lean{theorem3_maxdeg}~\cite{Tulsyan26}, which outputs $\HC(H)$ within
$\lfloor C2^{Cn\log_2(\Delta(H)+2)/\sqrt{\log_2(n+2)}}\rfloor$ steps for every graph $H$ on $n$
vertices, with \A's degree package \lean{exists_dPackage}, again after \lean{eulW}. For $H=G'$ we
have $3m$ vertices, maximum degree at most $\max(2,\Delta)$, and $\log_2(3m+2)\ge\log_2(m+2)$, so the
exponent is at most $C'm\log_2(\Delta+2)/\sqrt{\log_2(m+2)}$; the reading cost again gives
$(n+2)^C$. Graphs that are not codes are sent to output $0$ inside the machine, as in our directed
theorems.

\begin{remark}[Multigraphs; on paper]\label{rem:multi}
For a loopless multigraph $G$, subdividing every edge once gives a simple graph $G_1$ with $2m$ edges
and the same maximum degree when $\Delta\ge2$. Each new vertex has degree $2$ and a single
transition, so the transition systems of $G$ and $G_1$ correspond and $\EC(G_1)=\EC(G)$. So the
theorems extend to loopless multigraphs given as edge lists, with $m$ replaced by $2m$. Loops need a
convention about their two ends, and we do not treat them. None of this is formalised.
\end{remark}

% =====================================================================
\section{Numerical checks}\label{sec:numerics}
% =====================================================================

Before writing the proofs we computed, for twelve small graphs, three numbers that should agree, in a
Lean scratch file (\lean{lean/EthFollowup/Scratch/Euler/Check.lean}, evaluated with \texttt{\#eval},
Lean v4.34.1): Kotzig's count (transition systems with connected transition graph), $\HC(G')$ by
depth-first search on the arithmetic adjacency used by the word map, and the number of closed
Eulerian trails that start with edge $0$ in a fixed direction. These computations are separate from the
proof and are not machine-checked theorems. In a second scratch file (\lean{Check2.lean}) the kernel
evaluates the formal definitions themselves with \texttt{decide +kernel} (no \texttt{native\_decide}):
\lean{(eulerCircuitsAll G).card = 1} for the triangle and for $C_4$, and
\lean{(hamCycles (eulerSplitAll G)).card = 1} for the triangle.

\begin{table}[ht]
\centering\small
\begin{tabular}{@{}lrrrrr@{}}
\toprule
graph & $m$ & Kotzig & $\HC(G')$ & closed trails & $|E(G')|$\\
\midrule
$K_5$ & 10 & 132 & 132 & 132 & 50\\
octahedron $K_{2,2,2}$ & 12 & 372 & 372 & 372 & 60\\
$C_7^2$ (circulant, jumps $1,2$) & 14 & 1052 & 1052 & 1052 & 70\\
$K_4$ with two pendant triangles & 10 & 44 & 44 & 44 & 46\\
$K_{2,4}$ & 8 & 6 & 6 & 6 & 32\\
bowtie & 6 & 2 & 2 & 2 & 22\\
$C_6$ & 6 & 1 & 1 & 1 & 18\\
$K_4$ & 6 & 0 & 0 & 0 & 24\\
$K_5-e$ & 9 & 0 & 0 & 0 & 42\\
$K_{3,3}$ & 9 & 0 & 0 & 0 & 36\\
triangle & 3 & 1 & 1 & 1 & 9\\
$C_4$ & 4 & 1 & 1 & 1 & 12\\
\bottomrule
\end{tabular}
\caption{Checked numerically: three independent counts agree on every graph, and
$|E(G')|=2m+\sum_x\binom{d_x}{2}$ in every row. ``$K_4$ with two pendant triangles'' has the edges
$01,02,03,12,13,23,04,14,25,35$. $K_4$, $K_5-e$ and $K_{3,3}$ have odd degrees.}\label{tab:numerics}
\end{table}

\noindent A separate Python enumeration of all transition systems, by number of trails, gives
$243=132+96+15$ for $K_5$, $729=372+296+59+2$ for the octahedron, and $2187=1052+904+217+14$ for
$C_7^2$; the first terms agree with the table. The $K_5$ value agrees with OEIS
A007082~\cite{OEIS} (\cref{ex:k5}).

% =====================================================================
\section{The Lean development}\label{sec:formal}
% =====================================================================

In this section we list the modules and state the trusted base. The development is in the public
repository~\cite{TulsyanRepo26}. \Cref{tab:modules} lists the new modules.

\begin{table}[ht]
\centering\small
\begin{tabular}{@{}p{0.42\textwidth}p{0.43\textwidth}r@{}}
\toprule
module & content & lines\\
\midrule
\repo{lean/EthFollowup/ExactHC/EulerSplit.lean} & definitions of \cref{sec:defs}; bijection, size, arithmetic edges & 661\\
\repo{lean-bridge-anthropic/bridge/EthFollowup/ExactHC/EulerTheorems.lean} & word map \lean{eulW}, its polynomial time; the three theorems & 429\\
\midrule
total (proof) & & 1090\\
\midrule
\repo{lean/EthFollowup/Scratch/Euler/Check.lean}, \texttt{Check2.lean} & numerical checks, not part of the proof & 180, 98\\
\bottomrule
\end{tabular}
\caption{New modules, with line counts (\texttt{wc -l}). \texttt{EulerSplit} imports only
\texttt{HamiltonianParity} (431 lines) and Mathlib. \texttt{EulerTheorems} imports, unchanged, the
modules \texttt{DirectedTheorems} (790 lines), \texttt{DPackage} (548) and through them
\texttt{Theorem3} (1182), \texttt{Package} (1037), \texttt{Theorem2} (1764) and the rest of the verified
Hamiltonian-cycle development of~\cite{Tulsyan26}.}\label{tab:modules}
\end{table}

\paragraph{Toolchain.} The bridge builds on \A's Lean development at the commit pinned in our
repository (\texttt{anthropics/formal-math}, folder \texttt{eth}, commit \texttt{16a5ce4}, Lean
v4.33.1). The pure module \lean{EulerSplit} also compiles with Lean v4.34.1.

\paragraph{Trusted base and axioms.} \lean{\#print axioms} reports \lean{propext},
\lean{Classical.choice} and \lean{Quot.sound} for each of the three theorems. A fresh
\lean{leanchecker} replay of the dependency closure of the module
\lean{EthFollowup.ExactHC.EulerTheorems} passed on 4~November 2026. What must be trusted is the Lean
kernel; the Mathlib definitions used in the statements (\lean{SimpleGraph}, \lean{Connected},
\lean{Finset}, \lean{Sym2}, \lean{Nat.bits}, \lean{Nat.floor}, \lean{Real.logb}, \lean{Real.sqrt},
\lean{TM2OutputsInTime}); \A's two definitions \lean{GraphProblems.code} and
\lean{GraphProblems.Outputs}; and the definitions of \cref{sec:defs}. \A's development, our verified
Hamiltonian-cycle development~\cite{Tulsyan26} and Mathlib's proofs are checked by the kernel with ours;
they are part of the checked closure, not of the trusted base of definitions. But the soundness of the
whole rests on them building as a sound development (\cref{sec:limitations}).

\paragraph{What is not formalised.} The formal object is \lean{eulerCircuitsAll G}, Kotzig's
transition-system definition. We have not formalised that its size equals the number of closed
Eulerian walks of Mathlib (\lean{SimpleGraph.Walk.IsEulerian}) up to rotation and reversal; that
bridge (\cref{rem:kotzig}) is proved on paper only and checked numerically (\cref{tab:numerics},
column ``closed trails'').

% =====================================================================
\section{Related work}\label{sec:related}
% =====================================================================

We organise related work by idea: the complexity of counting Eulerian circuits, exact algorithms on
restricted classes, transitions and line graphs, and the counting engines with their formal proofs.

\paragraph{Complexity of counting.} For directed graphs the BEST theorem~\cite{BE51,ST41} gives a
polynomial-time count. For undirected graphs Brightwell and Winkler~\cite{BW05} proved
$\#$P-completeness, by a reduction from counting Eulerian orientations, and Ge and
\v{S}tefankovi\v{c}~\cite{GS10} extended it to planar $4$-regular graphs. Schrijver~\cite{Sch83}
bounds the number of Eulerian orientations, which is what the BEST-sum algorithm enumerates. In
contrast, we give no polynomial algorithm and no hardness result; we give an exact algorithm that is
subexponential in $m$, conditional on \A.

\paragraph{Exact algorithms on restricted classes.} Chebolu, Cryan and Martin~\cite{CCM12} count
Eulerian circuits in polynomial time on generalised series-parallel graphs, and Balaji, Datta and
Ganesan~\cite{BDG15} on graphs of bounded treewidth, even in parallel ($\#\mathrm{SAC}^1$). These are
unconditional and practical on their classes. In contrast, our bound holds for all graphs of bounded
degree, including expanders, but is conditional and galactic.

\paragraph{Transitions and line graphs.} Describing Eulerian circuits by transitions goes back to
Kotzig~\cite{Kot68}, who also studied Eulerian circuits with forbidden
transitions~\cite{Kot68b}. The bijection between Eulerian circuits and Hamiltonian cycles of a line
graph of a subdivision is classical, and Balaji, Datta and Ganesan~\cite{BDG15} count Euler tours by
modifying $G$ to a graph $G'$ with the same number of Euler tours and counting the Hamiltonian cycles of
$L(G')$ (their text does not fix $G'$ as the double subdivision). Closest to our work is that reduction; in contrast,
we use the bijection as a reduction to a Hamiltonian-cycle counter, and we prove it in Lean.

\paragraph{The counting engines and formal proofs.} OpenAI~\cite{OpenAI26} and \A~\cite{Anthropic26}
claim deterministic subexponential $\#$SAT algorithms, with Lean developments, under review. Our earlier
work~\cite{Tulsyan26} counts Hamiltonian cycles on top of \A's development; this note uses its
theorems as black boxes. We use Lean~4~\cite{Lean4} and Mathlib~\cite{Mathlib20}. In contrast to our
earlier work, the new mathematics here is small: a classical bijection and its size bounds, which we
state and check in full.

% =====================================================================
\section{Limitations}\label{sec:limitations}
% =====================================================================

\begin{itemize}
\item \emph{The result stands or falls with \A's development.} The proof rests on Anthropic's Lean
development of its claimed refutation of the Exponential Time Hypothesis, which is under review, and
on our Hamiltonian-cycle development built on it. Lean checks both as part of our proof, so we do not
need \A's printed theorems to match its formal ones; we do need the development, built from source at
the pinned commit, to be sound, and the toolchain to be sound.
\item \emph{The constants are astronomical.} They come from \A's levels and from our network
constructions, and none is optimised. The algorithm is of no practical use.
\item \emph{The formal object is Kotzig's definition.} The bridge to Mathlib's Eulerian walks is not
formalised (\cref{sec:formal}).
\item \emph{Simple graphs only.} Multigraphs reduce to simple graphs by subdivision
(\cref{rem:multi}), on paper only; loops are not treated.
\item \emph{Dense graphs are not covered.} \Cref{thm:rate} is subexponential in $m$ only when
$\log\Delta=o(\sqrt{\log m})$; \cref{thm:euler} needs bounded degree.
\item \emph{Variants are in progress and not claimed.} The bijection of \cref{lem:bij} holds for every
transition predicate $A$, which covers forbidden transitions, but the Turing-machine theorems for
forbidden transitions, for the minimum-cost Eulerian circuit with transition costs, and for outputting
a circuit are in progress, and we claim none of them here.
\item \emph{No human has checked the work.} The Lean code and this text were written with AI
assistance; no human expert has yet checked that \cref{sec:defs} lists everything a reader must
trust.
\end{itemize}

\section*{Acknowledgements}
This work was done with AI assistance (Claude Code, Anthropic): the reduction, the Lean development,
the numerical checks, the literature and citation checks and the text were produced in sessions with
Claude Code. The Lean theorems are checked by the Lean kernel; everything else awaits human expert
review. Responsibility for any errors rests with the author. This note continues the project
of~\cite{Tulsyan26}.

\begin{thebibliography}{99}

\bibitem{Anthropic26}
Claude (Anthropic), \emph{A $2^{o(n)}$ time algorithm for {3-SAT}}, preprint dated 26~October 2026.
Paper and Lean formalisation in the repository \url{https://github.com/anthropics/formal-math}, folder
\texttt{eth}; we use commit \texttt{16a5ce41304fcdaea4e0be7c132e1c361dc26da5}.

\bibitem{BDG15}
N.~Balaji, S.~Datta and V.~Ganesan, \emph{Counting {E}uler tours in undirected bounded treewidth
graphs}, 35th IARCS Annual Conference on Foundations of Software Technology and Theoretical Computer
Science (FSTTCS 2015), LIPIcs, vol.~45, 2015, pp.~246--260; arXiv:1510.04035.

\bibitem{BE51}
T.~van Aardenne-Ehrenfest and N.~G.~de Bruijn, \emph{Circuits and trees in oriented linear graphs},
Simon Stevin \textbf{28} (1951), 203--217.

\bibitem{BW05}
G.~Brightwell and P.~Winkler, \emph{Counting {E}ulerian circuits is $\#${P}-complete}, Proc. 7th
Workshop on Algorithm Engineering and Experiments and 2nd Workshop on Analytic Algorithmics and
Combinatorics (ALENEX/ANALCO 2005), SIAM, 2005, pp.~259--262; arXiv:cs/0405067.

\bibitem{CCM12}
P.~Chebolu, M.~Cryan and R.~Martin, \emph{Exact counting of {E}uler tours for generalized
series-parallel graphs}, J. Discrete Algorithms \textbf{10} (2012), 110--122.

\bibitem{GS10}
Q.~Ge and D.~\v{S}tefankovi\v{c}, \emph{The complexity of counting {E}ulerian tours in $4$-regular
graphs}, LATIN 2010: Theoretical Informatics, Lecture Notes in Comput. Sci., vol.~6034, Springer,
2010, pp.~638--649.

\bibitem{Kot68}
A.~Kotzig, \emph{Eulerian lines in finite $4$-valent graphs and their transformations}, Theory of
Graphs (Proc. Colloq., Tihany, 1966), Academic Press, New York, 1968, pp.~219--230.

\bibitem{Kot68b}
A.~Kotzig, \emph{Moves without forbidden transitions in a graph}, Mat. \v{C}asopis \textbf{18}
(1968), no.~1, 76--80.

\bibitem{Lean4}
L.~de Moura and S.~Ullrich, \emph{The {L}ean 4 theorem prover and programming language}, Automated
Deduction -- CADE 28, Lecture Notes in Comput. Sci., vol. 12699, Springer, 2021, pp.~625--635.

\bibitem{Mathlib20}
The mathlib Community, \emph{The {L}ean mathematical library}, Proc. 9th ACM SIGPLAN International
Conference on Certified Programs and Proofs (CPP 2020), 2020, pp.~367--381.

\bibitem{OEIS}
OEIS Foundation, \emph{The On-Line Encyclopedia of Integer Sequences}, sequence A007082,
\url{https://oeis.org/A007082}.

\bibitem{OpenAI26}
OpenAI, \emph{The {E}xponential {T}ime {H}ypothesis is false}, preprint dated 17~October 2026, revised
25~October 2026. Paper and Lean formalisation in the repository \url{https://github.com/openai/math}.

\bibitem{Sch83}
A.~Schrijver, \emph{Bounds on the number of {E}ulerian orientations}, Combinatorica \textbf{3}
(1983), no.~3--4, 375--380.

\bibitem{ST41}
C.~A.~B.~Smith and W.~T.~Tutte, \emph{On unicursal paths in a network of degree $4$}, Amer. Math.
Monthly \textbf{48} (1941), 233--237.

\bibitem{Tulsyan26}
V.~Tulsyan, \emph{Counting {H}amiltonian cycles in sparse graphs in subexponential time, assuming the
new $\#${SAT} algorithms}, Hexagon preprint hexagon:2610.00672v1, October 2026; source and Lean
development at \url{https://github.com/xennygrimmato/eth-hamiltonicity}.

\bibitem{TulsyanRepo26}
V.~Tulsyan, \emph{eth-hamiltonicity: {L}ean 4 formalisation and research notes}, public repository
\url{https://github.com/xennygrimmato/eth-hamiltonicity}, 2026; Lean sources as of commit \texttt{394740c}.

\end{thebibliography}

% =====================================================================
% CITATION CHECK (4 November 2026). Each entry was checked by web search or by fetching the source:
%  BW05   arXiv cs/0405067 (Brightwell, Winkler, submitted 18 May 2004, arXiv title "Note on counting
%         Eulerian circuits"; Thm 2.1: counting Eulerian circuits is #P-complete); published as
%         "Counting Eulerian circuits is #P-complete", ALENEX/ANALCO 2005 (researchr BrightwellW05;
%         LSE eprint 13350). Pages 259-262 from the standard citation; not seen on a publisher page.
%  GS10   Ge, Stefankovic, LATIN 2010, LNCS 6034, 638-649, DOI 10.1007/978-3-642-12200-2_55
%         (latintcs.org bib); #P-complete for planar 4-regular graphs; arXiv 1009.5019.
%  BDG15  Balaji, Datta, Ganesan, arXiv 1510.04035 (v2, Dec 2015), FSTTCS 2015, LIPIcs 45, 246-260.
%         Full text checked 4 Nov 2026 (pdftotext): they count Euler tours of G as Hamiltonian tours of L(G') for a
%         modified G' with tw(G') <= tw(G)+3; the text does not identify G' as the double subdivision, so the paper
%         attributes to them only the line-graph reduction, not our exact construction.
%         (Dagstuhl DROPS). v1 claimed GapL, retracted in v2; we cite only bounded treewidth and
%         #SAC^1. Line graph of the double subdivision: as stated in our EulerSplit.lean docstring,
%         taken from the planner's reading of the paper; not re-read in this pass.
%  CCM12  Chebolu, Cryan, Martin, J. Discrete Algorithms 10 (2012) 110-122,
%         DOI 10.1016/j.jda.2011.03.011 (DBLP journals/jda/CheboluCM12; arXiv 1005.3477).
%  BE51   van Aardenne-Ehrenfest, de Bruijn, Simon Stevin 28 (1951) 203-217 (Wikipedia BEST theorem,
%         MathWorld, Canad. J. Math. reference list agree).
%  ST41   Smith, Tutte, Amer. Math. Monthly 48 (1941) 233-237 (course reference list; consistent with
%         the BEST theorem literature). Volume/pages not seen on JSTOR.
%  Kot68  Kotzig, Theory of Graphs (Proc. Colloq. Tihany, Sept. 1966; eds. Erdos, Katona),
%         Academic Press 1968, 219-230 (Springer reference and a bibliography agree).
%  Kot68b Kotzig, Mat. Casopis 18 (1968), no. 1, 76-80 (EuDML 33972, fetched).
%  Sch83  Schrijver, Combinatorica 3 (1983), no. 3-4, 375-380 (CWI repository 10053); lower bound
%         (binom(d,d/2)/2^(d/2))^n for d-regular graphs, d even.
%  OEIS   A007082 fetched: K5 value 264 (reversed tours counted separately); "if we fix the Euler
%         tour to start with the edge 12, we get 132 Euler tours".
%  Lean4, Mathlib20, Anthropic26, OpenAI26, Tulsyan26: reused from the audited entries of the earlier
%         notes (research/citation-audit.md, 30 Oct 2026); Anthropic commit from
%         lean-bridge-anthropic/setup.sh (REV=16a5ce41...).
%  TulsyanRepo26: commit from `git log -1 --format=%h --
%         lean-bridge-anthropic/bridge/EthFollowup/ExactHC/EulerTheorems.lean` = 394740c.
% =====================================================================

\end{document}
